\chapter{Simulation Versus Live}\label{ch:rs:simulation-versus-live}

The research replay of a quoting strategy, the level-3 simulation of chapter~18 run on the same hours the strategy then traded, said it would lose
about \$50 an hour. Live, it lost \$491 an hour, and filled twice as many lots. Nothing was wrong with the replay's code; its fills were reproduced
in three languages. What went wrong is what every firm meets on its first live month: the simulator answered a slightly different question from the one
the market asked. This chapter builds the reconciliation (\texttt{firm.simlive}): match the live fills to the simulated ones, walk the simulated P\&L to
the live one an assumption at a time, locate the fills the simulator could not produce, and decide what to change so that the next estimate is closer. The
``live'' market is a stand-in that can be rerun: \texttt{firm.tape} with the strategy trading inside it as an agent, its orders real orders that other traders
see and trade against.

\section{Reconciling fills}

\begin{definition}[Sim-to-live reconciliation, shadow trading]\label{def:rs:simulation-versus-live:reconciliation}
\emph{Sim-to-live reconciliation}\index{sim-to-live reconciliation} is the systematic comparison of a strategy's live orders, fills and P\&L with what its
simulator produces on the same market data and the same signals. \emph{Shadow trading}\index{shadow trading} runs the simulator in parallel with live trading,
on the live feed, so that the comparison is available every day.
\end{definition}

The chapter's quoter keeps one lot on the best bid and one on the best ask within five lots of inventory. Live, inside \texttt{firm.tape}, it draws its own
latencies (market data 0.02 seconds plus an exponential of mean 0.02; order entry 0.03 plus an exponential of mean 0.02) and pays 0.1 tick a share. Its orders
are matched in time priority, count in the book that informed traders and liquidity providers react to, and are never cancelled by the background flow. The
research estimate is the replay of chapter~18 (fifo queue model) on the same hour without the quoter: the same seed and the same efficient-price path, as a
researcher would have run it the day before. Over six one-hour sessions the quoter filled 597 lots an hour live; the replay, on the live hour's messages without
the quoter's own and with the live latencies, 305.

\begin{definition}[Replay parity test]\label{def:rs:simulation-versus-live:parity}
A \emph{replay parity test}\index{replay parity test} replays a live session through the simulator, with the live orders' decisions and timings, and measures the
share of live fills the simulator reproduces (same side, price and time within a tolerance) and the share of simulated fills that happened live.
\end{definition}

\texttt{firm.simlive.match\_fills} pairs each live fill with a simulated fill of the same side and price within a second. Only 27\% of the live lots have a simulated
counterpart, and 54\% of the simulated lots happened live. A parity this low says the simulator is not a noisy version of live; it is a different model of it.

\section{Reconciling P\&L}

\begin{definition}[Divergence decomposition]\label{def:rs:simulation-versus-live:decomposition}
A \emph{divergence decomposition}\index{divergence decomposition} walks the simulated P\&L to the live P\&L through a sequence of intermediate estimates, each changing
one assumption to its live value, and attributes each step's change to that assumption.
\end{definition}

The order of the steps matters (changing latency first and the market second is not the same as the reverse), and each step must be a run of the simulator, not an
estimate. The chapter's sequence, averaged over the six sessions (\cref{fig:rs:simulation-versus-live:waterfall}):
\begin{enumerate}
\item the research replay: no latency, no fees, the hour without the quoter: $-\$50.5$ an hour;
\item with the live mean latencies: $-\$56.4$ (latency: $-\$5.9$);
\item on the market the quoter actually met, the live hour's messages without its own: $-\$112.8$ (the market: $-\$56.4$);
\item the live fills themselves, before fees: $-\$431.2$ (the fills: $-\$318.3$);
\item live, after fees: $-\$490.9$ (fees: $-\$59.7$).
\end{enumerate}

The market step compares two different realisations of the hour's order flow, with and without the quoter, and part of it is chance. A second realisation of each hour
without the quoter (the same efficient-price path, another order-flow seed) differs from the first by \$62.8 an hour on average: the market step is within its chance.
Latency and fees are small and known; the gap is in the fills.

\begin{omfigure}
\begin{tikzpicture}
\begin{axis}[width=11cm,height=5.8cm,ybar,bar width=20pt,ylabel={P\&L an hour (\$)},tick label style={font=\small},label style={font=\small},
  xtick={0,1,2,3,4},xticklabels={{research\\replay},{+ live\\latency},{+ market\\met},{+ live\\fills},{+ fees:\\live}},
  xticklabel style={align=center,font=\footnotesize},xmin=-0.6,xmax=4.6,ymin=-560,ymax=0,grid=major,grid style={gray!25},
  nodes near coords,nodes near coords style={font=\footnotesize,/pgf/number format/fixed,/pgf/number format/fixed zerofill,/pgf/number format/precision=1},
  table/col sep=comma]
\addplot[fill=omThm!45,draw=omThm] table[x=k,y=level]{figdata/research/19-simulation-versus-live/waterfall.csv};
\end{axis}
\end{tikzpicture}

\omcaption{From the research replay to live, an assumption at a time: the quoter's P\&L an hour, averaged over six simulated sessions, at each step of the
decomposition, each \omterm{def:rs:market-data-for-research:bar}{bar} the level after that step. Data: \texttt{rs\_simlive.decomposition}.}\label{fig:rs:simulation-versus-live:waterfall}
\end{omfigure}

\section{Locating the divergence}

Two mechanisms produce the live fills the replay cannot. They are visible only in the live messages, with the quoter's own orders in them.

\textbf{Trades that stopped at the quoter.} An aggressive order that meets the quoter's lot at the front of the queue is filled by it and may stop there. In the live
messages without the quoter's, that trade does not exist; a replay of them cannot fill the shadow with volume it never sees. 68\% of the live lots came from aggressive
orders that executed only against the quoter. The market after a firm trades is not the market without the firm.

\textbf{The last order at a price.} When the other orders at a price are cancelled as the price moves away, the quoter, whose cancellation is still travelling to the
exchange, is left alone at a stale price, and the next aggressor takes it. In the replay the level simply empties, and nothing trades there. 21\% of the live lots were
filled with no other order resting at their price, and their mark-out after ten seconds was $-0.85$ ticks, against $-0.11$ for the other live fills: the worst fills
of the day are the ones the simulator never shows.

\section{Keeping the simulator honest}

\begin{definition}[Simulator calibration, implementation shortfall, arrival price]\label{def:rs:simulation-versus-live:calibration}
\emph{Simulator calibration}\index{simulator calibration} adjusts a simulator's parameters (latency, queue model, fill probabilities) so that its output matches live
outcomes. The \emph{implementation shortfall}\index{implementation shortfall} of an order (Perold, 1988) is the difference between the P\&L of a paper portfolio traded
in full at the \omterm{def:rs:vectorised-backtests:timing}{decision time}'s price and the P\&L actually realised; the price at the \omterm{def:rs:vectorised-backtests:timing}{decision time}, or at the order's arrival in the market, is its \emph{arrival
price}\index{arrival price}.
\end{definition}

Calibration is the obvious response, and here it fails instructively. No latency makes the replay fill as many lots as live did: even at zero latency the replay on the
live messages fills 1\,898 lots over the six sessions against live's 3\,582. Switching to the front-of-queue model fills 8\,886, more than live, with a P\&L of $+\$2\,207$
against live's $-\$2\,587$. A simulator tuned to match the fill count would have tripled the error in P\&L. The divergence is structural, and the fix is structural too: a
reactive simulation, in which the strategy's orders are part of the market (as the agent here, and as the exchange simulator of Book~10), for strategies whose orders
are a noticeable share of the flow at their prices; and for the rest, a replay kept honest by the numbers this chapter produces every day.

The standing reconciliation is a report, produced daily from \omterm{def:rs:simulation-versus-live:reconciliation}{shadow trading}: parity of fills (both shares), the decomposition's steps, the mark-outs of live fills by
group (matched, live only, alone at the price), and the \omterm{def:rs:simulation-versus-live:calibration}{implementation shortfall} of every parent order against its \omterm{def:rs:simulation-versus-live:calibration}{arrival price}, split into execution, opportunity and
fees (\texttt{firm.simlive.implementation\_shortfall}). A simulator that has not been reconciled with live is a hypothesis; one that is reconciled every day is an
instrument.

\begin{omfigure}
\begin{tikzpicture}
\begin{axis}[width=11cm,height=5.4cm,ybar,bar width=8pt,xlabel={session},ylabel={P\&L (\$ an hour)},tick label style={font=\small},label style={font=\small},
  xtick={1,2,3,4,5,6},ymin=-1050,ymax=150,grid=major,grid style={gray!25},legend columns=3,
  legend style={font=\footnotesize,at={(0.5,-0.3)},anchor=north,draw=none,fill=none,/tikz/every even column/.append style={column sep=8pt}},
  table/col sep=comma]
\addplot[fill=omDef!55,draw=omDef] table[x=session,y=sim]{figdata/research/19-simulation-versus-live/sessions.csv};
\addplot[fill=gray!40,draw=gray] table[x=session,y=chance]{figdata/research/19-simulation-versus-live/sessions.csv};
\addplot[fill=omThm!55,draw=omThm] table[x=session,y=live]{figdata/research/19-simulation-versus-live/sessions.csv};
\legend{research replay,replay of another realisation,live}
\end{axis}
\end{tikzpicture}

\omcaption{The quoter's P\&L in each of six simulated sessions: the research replay, the replay of a second realisation of the same hour (with the live latencies), and live.
The live loss exceeds both in every session. Data: \texttt{rs\_simlive.decomposition}.}\label{fig:rs:simulation-versus-live:sessions}
\end{omfigure}

\section{Tutorial: ten times the simulated loss}\label{tut:rs:simulation-versus-live:tut}

\begin{tutorial}
\textbf{Goal.} Trade the quoter live inside \texttt{firm.tape}, replay the same hours, and reconcile: parity, decomposition, the location of the divergence, and a
calibration that fails. \textbf{End state:} \cref{fig:rs:simulation-versus-live:waterfall,fig:rs:simulation-versus-live:sessions}; the numbers of this chapter.
\begin{tutsteps}
\item \textbf{Match the fills}: greedy, in time order, by side and price within a tolerance.
\omcode{code/firm/simlive/firm_simlive.py}{31}{51}{Matching live fills to simulated ones.}\label{lst:rs:simulation-versus-live:match}
\item \textbf{The agent}: the quoter inside the market, with its own latencies.
\omcode{code/research/19-simulation-versus-live/python/rs_simlive.py}{36}{70}{The live quoter as a firm.tape agent.}\label{lst:rs:simulation-versus-live:agent}
\item \textbf{Run} \texttt{decomposition()}, \texttt{steps()}, \texttt{vanished\_volume()}, \texttt{alone\_at\_price()}, \texttt{effective\_latency()},
\texttt{model\_calibration()} and \texttt{fig\_simlive.py}.
\end{tutsteps}
\textbf{What to change next.} Make the quoter cancel as soon as the other orders at its price fall below a threshold, and watch the ``alone'' fills and the gap shrink;
restore the executions against the quoter in the replay's input as executions against an anonymous order, and measure the parity again.
\end{tutorial}

\section{Build: the reconciliation toolkit}\label{bld:rs:simulation-versus-live:simlive}

\begin{build}
\bfield{Purpose} The daily comparison of every live strategy with its simulator, and the evidence for changing the simulator.
\bfield{Interface} \texttt{FILL}, \texttt{as\_fills}, \texttt{match\_fills(live, sim, tol)}, \texttt{parity}, \texttt{fill\_pnl(fills, mark)}, \texttt{waterfall(steps)},
\texttt{calibrate(grid, metric, target)}, \texttt{implementation\_shortfall(side, qty, decision\_px, fills, end\_px, fee\_per\_unit)}; \texttt{firm.tape.simulate(cfg,
agent)} for the reactive stand-in.
\bfield{Rules} Every step of a decomposition is a run, in a stated order; parity is reported both ways; calibration is judged on P\&L and mark-outs, not on fill counts alone.
\bfield{Acceptance tests} \texttt{code/firm/simlive/tests/}: matching with partial quantities and a tolerance; parity both ways, and perfect on identical fills; per-fill P\&L
summing to the marked total; the waterfall's steps; calibration on a grid; the \omterm{def:rs:simulation-versus-live:calibration}{implementation shortfall} of a buy and a sell by hand.
\bfield{Stretch} Matching by order identifier with the exchange's acknowledgements; decomposition by time of day; recalibration of a probabilistic queue model on live fills.
\end{build}

\begin{omsources}
\item A.~F.~Perold, ``The implementation shortfall: paper versus reality'', \emph{Journal of Portfolio Management} 14(3), 1988.
\end{omsources}

\section{Exercises}

\begin{exercise}[$\star$]\label{exo:rs:simulation-versus-live:1}
Live fills: buy 100 at 50.00 at 10.2 s, sell 100 at 50.02 at 14.0 s. Simulated fills: buy 100 at 50.00 at 10.9 s, buy 100 at 49.99 at 12.0 s. With a one-second
tolerance, what are the two parity shares?
\end{exercise}

\begin{exercise}[$\star$]\label{exo:rs:simulation-versus-live:2}
A decision to buy 10\,000 shares is taken at 20.00. 6\,000 fill at an average of 20.04 and the rest is cancelled; the price ends at 20.30. Fees are 0.2 cents a share.
Compute the \omterm{def:rs:simulation-versus-live:calibration}{implementation shortfall} and its three parts.
\end{exercise}

\begin{exercise}[$\star$]\label{exo:rs:simulation-versus-live:3}
Why must each step of a \omterm{def:rs:simulation-versus-live:decomposition}{divergence decomposition} be a run of the simulator rather than an estimate?
\end{exercise}

\begin{exercise}[$\star\star$]\label{exo:rs:simulation-versus-live:4}
Why can the market step of the decomposition be within chance even when the strategy's impact is real? How would you separate the two?
\end{exercise}

\begin{exercise}[$\star\star$]\label{exo:rs:simulation-versus-live:5}
Explain why a replay of the live session without the strategy's own messages undercounts its fills, and what share of the chapter's live lots it cannot see.
\end{exercise}

\begin{exercise}[$\star\star$]\label{exo:rs:simulation-versus-live:6}
Why are fills that happen when the strategy is the last order at its price worse than the others?
\end{exercise}

\begin{exercise}[$\star\star\star$]\label{exo:rs:simulation-versus-live:7}
\emph{Coding.} With \texttt{rs\_simlive.effective\_latency}, confirm that no latency closes the fill gap. Which of the chapter's two mechanisms would a latency change
affect, and why not the other?
\end{exercise}

\begin{exercise}[$\star\star\star$]\label{exo:rs:simulation-versus-live:8}
\emph{Find the flaw.} ``Live fills were double the simulator's, so we switched the simulator to optimistic fills; now the fill counts match and the \omterm{def:rs:vectorised-backtests:backtest}{backtest} is
calibrated.''
\end{exercise}

\section{Problem: Ten Times the Simulated Loss}

\begin{problem}[{Weekend problem --- a live strategy against its simulator}]\label{pb:rs:simulation-versus-live:1}
The chapter's quoter, live inside \texttt{firm.tape} for six one-hour sessions, and its level-3 replays.

\medskip\noindent\textbf{Part I --- Fills.}
\begin{enumerate}
\item How many lots an hour did the quoter fill live, and how many did the replay fill on the live messages without its own?
\item What are the two parity shares?
\item What does a parity this low say about the simulator?
\item Which live latencies did the quoter have, and what does the research replay assume?
\end{enumerate}

\medskip\noindent\textbf{Part II --- P\&L.}
\begin{enumerate}[resume]
\item Give the P\&L an hour at each step of the decomposition.
\item How large is each step?
\item How much of the market step can chance explain?
\item Which steps are known in advance, and which had to be measured?
\end{enumerate}

\medskip\noindent\textbf{Part III --- The divergence.}
\begin{enumerate}[resume]
\item What share of the live lots came from aggressive orders that executed only against the quoter, and why can the replay not see them?
\item What share of the live lots were filled with no other order at their price, and what were their mark-outs against the others'?
\item Why does the quoter end up alone at a price?
\item What would reduce those fills?
\end{enumerate}

\medskip\noindent\textbf{Part IV --- The verdict.}
\begin{enumerate}[resume]
\item State the \emph{named result}: the decomposition of the gap between live and simulated P\&L into latency, the market, the fills and fees.
\item What does calibrating latency achieve?
\item What does switching to the front-of-queue model do to fills and P\&L?
\item What is the structural fix?
\item What should the daily reconciliation report contain?
\item How would you use the \omterm{def:rs:simulation-versus-live:calibration}{implementation shortfall} for this strategy's inventory unwinds?
\item What changes if the strategy trades 1\% of the volume at its prices instead of 10\%?
\item In one sentence: what makes a simulator an instrument rather than a hypothesis?
\end{enumerate}
\end{problem}

\section{Interview questions}

\begin{interviewq}[$\star$ \iqroles{researcher, trader}]\label{iq:rs:simulation-versus-live:1}
Your strategy makes half in production what it made in simulation. Where do you look first?
\end{interviewq}

\begin{interviewq}[$\star\star$ \iqroles{researcher}]\label{iq:rs:simulation-versus-live:2}
What is \omterm{def:rs:simulation-versus-live:calibration}{implementation shortfall}, and how do you decompose it?
\end{interviewq}

\begin{interviewq}[$\star\star$ \iqroles{developer, researcher}]\label{iq:rs:simulation-versus-live:3}
Describe a \omterm{def:rs:simulation-versus-live:parity}{replay parity test} and what it catches.
\end{interviewq}

\begin{interviewq}[$\star\star$ \iqroles{researcher, trader}]\label{iq:rs:simulation-versus-live:4}
Why does a \omterm{def:rs:vectorised-backtests:backtest}{backtest} on market data recorded while you were trading misstate your fills?
\end{interviewq}

\begin{interviewq}[$\star\star$ \iqroles{researcher}]\label{iq:rs:simulation-versus-live:5}
How would you calibrate a fill simulator on live fills without overfitting it?
\end{interviewq}

\begin{interviewq}[$\star\star\star$ \iqroles{researcher, developer}]\label{iq:rs:simulation-versus-live:6}
Design a daily \omterm{def:rs:simulation-versus-live:reconciliation}{sim-to-live reconciliation} for a firm with fifty strategies.
\end{interviewq}
